\section{Conclusion}

We began expecting a precision hierarchy --- that a more formally rigorous verifier would describe a defect more precisely and so produce better repair. Four categories, one backbone, 421 problems, and a Spearman correlation of $-1.0$ later, that expectation is not weakly supported but exactly inverted. The verifier that says the most fixes the least, and the verifier that costs the least returns the most correctness per second.

\subsection{Summary}

Holding GPT-4o-mini, the problem set, the prompt template, and a three-iteration budget fixed while varying only the feedback category, we found that the four categories produce feedback volumes separated by four orders of magnitude (Kruskal-Wallis $\varepsilon^2 = 0.88$), which makes the comparison meaningful; that iteration-1 repair success falls monotonically as feedback volume rises, from mypy's 6.35\% down to Pyright's 4.76\%, giving $\rho = -1.0$; that normalizing correctness by wall-clock overhead reverses the ranking a second time, with static analysis at a 6.64 efficiency ratio against execution monitoring's 0.41 despite execution monitoring delivering twice the absolute pass@1 improvement; and that SMT-guided repair is not reachable at this model tier at all, with usable Z3 constraints extractable for 6.3\% of problems against the 40\% we assumed.

Our account of the inversion is that repair is bounded by the model's ability to extract one actionable edit, not by the feedback's information content. A short traceback names the fix; a 24,000-character diagnostic tree contains it somewhere.

Two of these results came from experiments that failed pre-registered gates. We report them as the contribution because a refuted prediction that reverses a field-wide assumption is worth more than a confirmed one that restates it.

\subsection{Future Directions}

The most urgent work separates the two explanations we cannot currently distinguish. Pyright may underperform because verbose feedback is intrinsically harder to act on, or because our 4,000-character truncation discarded the decisive diagnostic, or because JSON rule names like \texttt{reportAttributeAccessIssue} are less grounded for a language model than natural-language error text. Comparing raw Pyright JSON against a one-sentence model-generated summary of the same diagnostics, on matched problems, would settle it: if the summary matches execution monitoring's repair rate while the raw JSON does not, the problem is presentation, and the fix is a summarization step costing one extra call.

Z3's position at the top of the inverse ordering needs the same treatment from the other direction. Matching the 8 Z3-feasible problems against the 118 infeasible ones on baseline difficulty, then comparing all four categories within the matched strata, would show whether Z3's apparent lead is a feedback effect or selection bias.

Our strongest assumption is the one we could not test: that GPT-4o-mini's behavior generalizes. The mechanism we propose predicts that a model better at parsing structured text should extract Pyright's decisive diagnostic more reliably, weakening or reversing the correlation. Running the three-category comparison on GPT-4o and Claude 3.5 Sonnet over the same 126 failing solutions would test the mechanism and the practical recommendation at once, since a sign reversal would invert our deployment advice rather than merely qualify it.

Beyond that, our bug distribution suggests a strategy we did not test. Logic errors account for 65.9\% of failures and type errors for 24.6\%, and our classifier separates them at 82\% agreement --- accurately enough to route. Sending predicted type errors to static analysis and logic errors to execution monitoring might beat any uniform assignment, and would cost roughly what our repair experiment cost to check. Extending the comparison to repository-level repair is the larger prize: in real codebases, where import errors and type mismatches are common and executing the test suite is expensive, static analysis should be worth relatively more than it is on function-level benchmarks.

\subsection{Closing}

The verify stage of a repair loop is usually designed around what the tool can prove. Our results suggest it should be designed around what the model can act on. Those turn out to be different objectives, and on current models they point in opposite directions.
